\documentclass[12pt, a4paper]{article}

\usepackage{amsmath, amssymb, amsthm}
\usepackage{geometry}
\usepackage{hyperref}
\usepackage{titlesec}
\usepackage{parskip}
\usepackage{microtype}
\usepackage{lmodern}
\usepackage{xcolor}
\usepackage{booktabs}
\usepackage{enumitem}
\usepackage{mathtools}

\geometry{margin=1.2in}

\definecolor{deepblue}{RGB}{20, 40, 100}
\definecolor{muted}{RGB}{90, 90, 90}

\hypersetup{colorlinks=true, linkcolor=deepblue,
            urlcolor=deepblue, citecolor=deepblue}

\titleformat{\section}{\large\bfseries\color{deepblue}}{{\thesection}}{1em}{}
\titleformat{\subsection}{\normalsize\bfseries\color{deepblue}}{{\thesubsection}}{1em}{}
\titleformat{\subsubsection}{\normalsize\itshape\color{deepblue}}{{\thesubsubsection}}{1em}{}

\theoremstyle{plain}
\newtheorem{theorem}{Theorem}[section]
\newtheorem{proposition}[theorem]{Proposition}
\newtheorem{lemma}[theorem]{Lemma}
\newtheorem{corollary}[theorem]{Corollary}
\theoremstyle{definition}
\newtheorem{definition}[theorem]{Definition}
\newtheorem{remark}[theorem]{Remark}
\newtheorem{example}[theorem]{Example}

% ---- notation ----
\newcommand{\R}{\mathbb{R}}
\newcommand{\C}{\mathbb{C}}
\newcommand{\Z}{\mathbb{Z}}
\newcommand{\N}{\mathbb{N}}
\newcommand{\T}{\mathbb{T}}
\newcommand{\Action}{\mathcal{A}}
\newcommand{\Lag}{\mathcal{L}}
\newcommand{\Fix}{\mathrm{Fix}}
\newcommand{\bfq}{\mathbf{q}}
\newcommand{\bfp}{\mathbf{p}}
\newcommand{\bfa}{\mathbf{a}}
\newcommand{\bfb}{\mathbf{b}}
\newcommand{\bfz}{\mathbf{z}}
\newcommand{\SO}{\mathrm{SO}}
\newcommand{\On}{\mathrm{O}}
\newcommand{\SU}{\mathrm{SU}}
\newcommand{\eps}{\varepsilon}
\newcommand{\abs}[1]{\left|#1\right|}
\newcommand{\norm}[1]{\left\|#1\right\|}
\newcommand{\ip}[2]{\langle #1,\, #2\rangle}

\begin{document}

% ---------- TITLE ----------
\begin{center}
  {\LARGE \textbf{Three-Dimensional Choreographic Orbits:\\[4pt]
  Mathematical Framework}}\\[1.6em]
  {\large \textcolor{muted}{MathFrame — v0.1
    \quad|\quad Companion to RDD v0.1 and CompSpec v0.1}}\\[0.6em]
  {\normalsize \textcolor{muted}{e5bed97b3349 \quad\textbullet\quad March 2026}}
\end{center}

\vspace{1em}
\noindent\rule{\textwidth}{0.4pt}

\begin{abstract}
\noindent
This document develops the mathematical foundations underlying the
three-dimensional choreographic orbit classification programme.  We
establish the variational framework, derive the Fourier-mode constraints
imposed by each relevant three-dimensional point group via representation
theory, state the symmetric-criticality existence theorem, develop the
knot-theoretic characterisation of the spacetime trace, and present the
Floquet--Krein stability theory for periodic Hamiltonian orbits.  These
results inform the research design and computational specification;
they will also constitute the theory section of the eventual paper.
\end{abstract}

\noindent\rule{\textwidth}{0.4pt}
\tableofcontents
\noindent\rule{\textwidth}{0.4pt}

% ============================================================
\section{Setup and Notation}
% ============================================================

\subsection{Units and normalisation}

Throughout we work in units where the gravitational constant $G = 1$,
all masses $m = 1$, and the choreographic period $T = 1$.  Distances
are measured in units such that the characteristic action
$\Action_0 = m^2 G / T^{1/3} = 1$.  Physical results are recovered
by dimensional analysis.

\subsection{Configuration space}

The configuration space of $N$ bodies in $\R^3$ is
\[
  X_N \;=\; (\R^3)^N \setminus \Delta,
  \qquad
  \Delta = \bigl\{\bfq \in (\R^3)^N : \bfq_i = \bfq_j
  \text{ for some } i \neq j\bigr\},
\]
with $\Delta$ the collision set.  The loop space we minimise over is
\[
  \Lambda_N \;=\; H^1(S^1,\, (\R^3)^N \setminus \Delta),
\]
the Sobolev space of absolutely continuous closed paths avoiding
collisions.  The $H^1$ norm ensures the kinetic-energy term is finite
and that $\Lambda_N$ carries a Hilbert manifold structure.

\subsection{The Newtonian action}

The Lagrangian action for one period is
\begin{equation}\label{eq:action}
  \Action[\bfq_1,\ldots,\bfq_N]
  \;=\;
  \int_0^1
  \left[
    \frac{1}{2}\sum_{i=1}^N \abs{\dot\bfq_i}^2
    \;+\;
    \sum_{i < j} \frac{1}{\abs{\bfq_i - \bfq_j}}
  \right] dt.
\end{equation}
Critical points of $\Action$ on $\Lambda_N$ are $T=1$ periodic solutions
of Newton's equations $\ddot\bfq_i = \sum_{j\neq i}(\bfq_j -
\bfq_i)/\abs{\bfq_j-\bfq_i}^3$.

% ============================================================
\section{The Choreographic Subspace}
% ============================================================

\subsection{Definition and reduction}

\begin{definition}[Choreographic orbit]\label{def:choreo}
A path $(\bfq_1,\ldots,\bfq_N) \in \Lambda_N$ is \emph{choreographic}
if there exists a single closed curve $\gamma \in H^1(S^1, \R^3)$ such
that
\[
  \bfq_i(t) = \gamma\!\left(t - \tfrac{i-1}{N}\right),
  \qquad i = 1,\ldots,N,
\]
where time is taken modulo 1.  The bodies follow the same curve in
sequence, each offset by phase $1/N$.
\end{definition}

Substituting Definition~\ref{def:choreo} into \eqref{eq:action} and
using $T$-periodicity collapses the functional to a single-curve action:
\begin{equation}\label{eq:single}
  \Action[\gamma]
  \;=\;
  \frac{1}{2}\int_0^1 \abs{\dot\gamma}^2\,dt
  \;+\;
  \frac{1}{N}\sum_{k=1}^{N-1}
  \int_0^1 \frac{dt}{\abs{\gamma(t) - \gamma(t + k/N)}}.
\end{equation}
The choreographic subspace $\mathcal{C}_N \subset \Lambda_N$ is the
image of the embedding $\gamma \mapsto (\bfq_1,\ldots,\bfq_N)$; it is
a closed submanifold of $\Lambda_N$.  Critical points of
$\Action|_{\mathcal{C}_N}$ correspond to genuine $N$-body solutions.

\subsection{Fourier representation}

Since $\gamma \in H^1(S^1,\R^3)$, it admits a Fourier expansion
\begin{equation}\label{eq:fourier}
  \gamma(t)
  = \sum_{k=1}^{\infty}
    \left[\bfa_k \cos(2\pi k t) + \bfb_k \sin(2\pi k t)\right],
  \qquad \bfa_k, \bfb_k \in \R^3,
\end{equation}
where the zero mode $\bfa_0 = 0$ by the centre-of-mass constraint
$\sum_i \bfq_i \equiv 0$.  The $H^1$ condition reads
$\sum_{k=1}^\infty k^2(\abs{\bfa_k}^2 + \abs{\bfb_k}^2) < \infty$.

It is convenient to pass to complex amplitudes.  Write
$\gamma(t) = \mathrm{Re}\sum_{k=1}^\infty \bfz_k\,e^{2\pi ikt}$, so
$\bfz_k = \bfa_k - i\bfb_k \in \C^3$.  The phase-shift
$\gamma(t) \mapsto \gamma(t + s)$ acts by $\bfz_k \mapsto
e^{2\pi iks}\bfz_k$; the spatial transformation $R \in \On(3)$ acts by
$\bfz_k \mapsto R\bfz_k$ (applied componentwise in $\C^3$ via the
complexification of $R$).

% ============================================================
\section{Symmetry Groups and Equivariant Critical Points}
% ============================================================

\subsection{The symmetry group of a choreography}

The symmetry group of a choreographic orbit is a subgroup of
\[
  G \;\subset\; \On(3) \times D_N,
\]
where $D_N = \Z/N\Z \rtimes \Z/2\Z$ is the dihedral group of order $2N$.
The $\Z/N\Z$ factor permutes the bodies cyclically; the $\Z/2\Z$ factor
is time-reversal $t \mapsto -t$, which for an equal-mass choreography
acts as a spatial isometry composed with reversal of the traversal
direction.  A generator of the symmetry group is a triple
$(R, \sigma, \epsilon)$ with $R \in \On(3)$, $\sigma \in \Z/N\Z$,
$\epsilon \in \{+1,-1\}$, acting on $\gamma$ by
\begin{equation}\label{eq:gaction}
  (R,\sigma,\epsilon)\cdot\gamma(t)
  \;=\;
  R\,\gamma\!\left(\epsilon t - \tfrac{\sigma}{N}\right).
\end{equation}
Most groups arising in practice have $\epsilon = +1$ (no time-reversal);
we specialise to this case initially and note where $\epsilon = -1$ is
required.

\subsection{Induced representation on Fourier modes}

Applying \eqref{eq:gaction} to the Fourier series \eqref{eq:fourier}
with $\epsilon = +1$ gives, for mode $k$:
\begin{equation}\label{eq:modeaction}
  \bfz_k \;\longmapsto\; e^{-2\pi ik\sigma/N}\,R\,\bfz_k.
\end{equation}
For $\gamma$ to be a fixed point of $(R,\sigma,+1)$, each mode must
satisfy
\begin{equation}\label{eq:fixcond}
  \bfz_k \;=\; e^{-2\pi ik\sigma/N}\,R\,\bfz_k,
\end{equation}
i.e.\ $\bfz_k$ lies in the $e^{2\pi ik\sigma/N}$-eigenspace of $R$
(in the complexification of $\R^3$).  This is the key equation from
which all symmetry constraints follow.

\subsection{The Principle of Symmetric Criticality}

\begin{theorem}[Palais, 1979]\label{thm:psc}
Let a compact Lie group $G$ act smoothly and isometrically on a
Hilbert manifold $\mathcal{M}$, and let $f : \mathcal{M} \to \R$ be a
$G$-invariant $C^1$ functional.  Then every critical point of
$f|_{\Fix(G)}$ is a critical point of $f$ on all of $\mathcal{M}$.
\end{theorem}

\begin{corollary}
Every local minimiser of $\Action|_{\Fix(G) \cap \mathcal{C}_N}$
is a genuine $N$-body choreographic solution.
\end{corollary}

The theorem applies here because $\On(3)$ is compact and acts
isometrically on $H^1(S^1,\R^3)$ with respect to the natural $L^2$
inner product, and $\Action$ is manifestly $\On(3)$-invariant.  The
Reynolds projector $P_G = |G|^{-1}\sum_{g\in G}\rho(g)$ is the
orthogonal projection onto $\Fix(G)$.

% ============================================================
\section{Point-Group Constraints: Fourier Mode Analysis}
% ============================================================

We now apply \eqref{eq:fixcond} systematically to each point-group
family.  For each group we state: the generator(s), the surviving
modes, and the dimension count.  We use the Sch\"{o}nflies notation
throughout.

\subsection{Cyclic groups $C_n$}

The group $C_n$ is generated by $(R_n, 1, +1)$, where $R_n \in
\SO(3)$ is rotation by $2\pi/n$ about the $z$-axis and $\sigma = 1$.
The complexified eigenvalues of $R_n$ are $\{e^{2\pi i/n}, e^{-2\pi
i/n}, 1\}$ for the three components $(z_+, z_-, z_0)$ in the
helical basis, where $z_\pm = x \pm iy$ and $z_0 = z$.

Equation \eqref{eq:fixcond} for this generator becomes:
\[
  \bfz_k \;=\; e^{-2\pi ik/n}\,R_n\,\bfz_k.
\]
Decomposing by helicity:
\begin{itemize}[leftmargin=1.8em]
  \item \textbf{$z_+$-component} (eigenvalue $e^{2\pi i/n}$):
        condition is $e^{2\pi i/n} e^{-2\pi ik/n} = 1$,
        i.e.\ $n \mid (1-k)$, i.e.\ $k \equiv 1 \pmod{n}$.
  \item \textbf{$z_-$-component} (eigenvalue $e^{-2\pi i/n}$):
        condition is $e^{-2\pi i/n} e^{-2\pi ik/n} = 1$,
        i.e.\ $n \mid (1+k)$, i.e.\ $k \equiv n-1 \pmod{n}$.
        Since $z_- = \overline{z_+}$, this is automatically satisfied
        when the $z_+$ condition holds.
  \item \textbf{$z_0$-component} ($z$-axis, eigenvalue $1$):
        condition is $e^{-2\pi ik/n} = 1$,
        i.e.\ $n \mid k$, i.e.\ $k \equiv 0 \pmod{n}$.
\end{itemize}

\begin{proposition}[$C_n$ mode selection]\label{prop:cn}
For an $N=n$ body choreography with $C_n$ symmetry, the surviving
Fourier modes are:
\[
  xy\text{-plane}: \quad k \;\in\; \{1, n+1, 2n+1, \ldots\},
  \qquad
  z\text{-axis}: \quad k \;\in\; \{n, 2n, 3n, \ldots\}.
\]
For a truncation at $K$ modes, the dimension of $\Fix(C_n)$ is
$\dim_\R = 4\lfloor K/n \rfloor + 2\lfloor (K-1)/n + 1\rfloor
\approx 6K/n$.
\end{proposition}

\begin{example}[$C_3$, $N=3$]
Surviving modes: $xy$ at $k = 1, 4, 7, 10, \ldots$;
$z$ at $k = 3, 6, 9, \ldots$.  At truncation $K = 12$:
4 xy-modes ($k=1,4,7,10$) contributing 8 real d.o.f.,
plus 4 z-modes ($k=3,6,9,12$) contributing 4 real d.o.f.
Total: 12 real parameters vs.\ the unconstrained 36. Reduction factor: 3.
\end{example}

\subsection{Dihedral groups $D_n$, $D_{nh}$, $D_{nd}$}

$D_n$ adds a $C_2$ rotation about an axis perpendicular to $z$, e.g.\
$R_{2x}$ (rotation by $\pi$ about $x$), with body permutation $\sigma
= 0$ (no phase shift).  In helicity components:
$R_{2x} : z_+ \mapsto -\overline{z_-}$, $z_0 \mapsto -z_0$.
The additional constraint on surviving $C_n$ modes is:

\begin{itemize}[leftmargin=1.8em]
  \item \textbf{$z_+$-component}: $R_{2x}$ maps the $k\equiv 1$
        modes to $k \equiv n-1$ modes.  Compatibility requires
        $\bfz_k^+ = -\overline{\bfz_k^-}$, fixing the imaginary part
        of $z_+$ in terms of the real part.  Net effect: one real
        parameter per surviving $xy$ mode instead of two.
  \item \textbf{$z$-component}: $R_{2x}$ maps $z_0 \mapsto -z_0$,
        so $k\equiv 0$ z-modes are forced to zero.
\end{itemize}

\begin{proposition}[$D_n$ mode selection]
For $D_n$ symmetry with $N=n$, the surviving modes are:
$xy$-plane modes at $k \equiv 1 \pmod{n}$ with \emph{one} real
d.o.f.\ each; $z$-axis modes vanish entirely.  The orbit is
therefore confined to a plane: $D_n$ choreographies with $N=n$
are generically planar.
\end{proposition}

This result is significant: it means that searching for non-planar
choreographies requires symmetry groups where the $z$-modes are
\emph{not} killed.  The $C_n$ and $C_{nv}$ families are the
natural targets for genuinely 3D orbits.

\textbf{$D_{nh}$ and $D_{nd}$:} Adding a horizontal reflection
$\sigma_h$ (for $D_{nh}$) maps $z_0 \to -z_0$, again
eliminating $z$-modes and enforcing planarity.  $D_{nd}$ adds
an $S_{2n}$ improper rotation; the analysis is similar and also
yields planar orbits for generic $N$.

\subsection{Cyclic groups with vertical reflection: $C_{nv}$}

$C_{nv}$ adds a vertical reflection $\sigma_v$ through the $xz$-plane,
which acts as $(x,y,z) \mapsto (x,-y,z)$, i.e.\ $z_+ \mapsto
\overline{z_+}$, with $\sigma=0$ (no phase shift).  The $C_n$ mode
selection survives; the $\sigma_v$ constraint additionally requires
$\bfz_k^{xy}$ to be real in the helicity basis, fixing
$\mathrm{Im}(\bfz_k^{xy}) = 0$.  For the $z$-modes, $\sigma_v$ acts
trivially, leaving them free.

\begin{proposition}[$C_{nv}$ mode selection]\label{prop:cnv}
For $C_{nv}$ with $N=n$:
$xy$-plane modes at $k \equiv 1 \pmod{n}$ each contribute one
real d.o.f.\ (not two); $z$-axis modes at $k \equiv 0 \pmod{n}$
each contribute two real d.o.f.\ (both $\bfa_k^z$ and $\bfb_k^z$ are free).
The orbits are generically \emph{non-planar}.
\end{proposition}

The $C_{nv}$ family therefore provides the simplest non-planar
choreographic orbits with non-trivial symmetry.  The $C_{3v}$ case
with $N=3$ is our primary Phase\,1 target.

\subsection{Tetrahedral group $T$ and $T_d$}

The tetrahedral group $T \cong A_4$ has order 12, generated by a
$C_3$ rotation about the $(1,1,1)$-axis and a $C_2$ rotation about
the $z$-axis.  For an $N$-body choreography with $T$ symmetry, $N$
must be divisible by the order of the $C_3$ element in the body-permutation
representation; the smallest natural choice is $N=12$ (one body per
element of $T$).

Applying \eqref{eq:fixcond} to all 12 generators of $T$ acting
simultaneously constrains the surviving modes.  The key result is:

\begin{proposition}[$T$ mode selection, $N=12$]\label{prop:T}
For $N=12$ with full tetrahedral symmetry $T$, the first
surviving mode index is $k=1$ (which carries a 3-dimensional
real space of coefficients transforming as the standard
representation of $T$).  Subsequent survivors occur at
$k \equiv 1 \pmod{12}$ for the $xy$-components and
$k \equiv 0 \pmod{12}$ for the $z$-component (the $A_1$
representation).  At truncation $K=48$ this gives approximately
$48/12 \times 6 \approx 24$ real parameters.
\end{proposition}

$T_d$ adds improper $S_4$ axes.  The additional constraint selects
among representations of $T$ and may further eliminate $z$-modes,
potentially restricting to planar orbits; this requires case-by-case
analysis.

\subsection{Octahedral group $O$ and $O_h$}

The octahedral group $O \cong S_4$ has order 24.  For $N=24$
(one body per group element), the analysis parallels the tetrahedral
case.  The mode constraints are sparser: surviving modes occur at
$k \equiv 1 \pmod{24}$ for the $xy$-components, and the first
non-trivial $z$-mode is at $k = 24$.  At truncation $K=96$:
approximately $96/24 \times 6 = 24$ real parameters — the same
dimensionality as the tetrahedral case, despite the larger group,
because the higher symmetry compensates for the larger $N$.

$O_h = O \times \Z/2\Z$ (adding inversion) tends to force additional
pairings between modes and can eliminate the $z$-modes entirely
depending on $N$.

\subsection{Icosahedral group $I$ and $I_h$}

$I \cong A_5$ has order 60, making it the largest finite subgroup of
$\SO(3)$.  For $N = 60$, the mode constraints are:
$xy$-components survive at $k \equiv 1 \pmod{60}$ and $z$-components
at $k \equiv 0 \pmod{60}$.  At truncation $K = 240$: $240/60 \times
6 = 24$ real parameters — the same order as tetrahedral and
octahedral cases.  The key open question is whether $N=12$ (a proper
subgroup arrangement) admits icosahedral-class choreographies.

\begin{remark}
The dimension count suggests that the search is no harder for
icosahedral than tetrahedral symmetry at fixed truncation ratio $K/N$
— the larger group buys just as much reduction as the larger $N$
costs.  This is the dimensional miracle that makes the icosahedral
case computationally tractable.
\end{remark}

\subsection{Summary table}

\begin{center}
\begin{tabular}{lllllc}
\toprule
\textbf{Group} & \textbf{Order} & \textbf{Natural $N$}
  & \textbf{$xy$ modes} & \textbf{$z$ modes}
  & \textbf{3D?} \\
\midrule
$C_n$   & $n$  & $n$  & $k \equiv 1 \pmod{n}$ & $k \equiv 0\pmod{n}$ & Yes \\
$C_{nv}$& $2n$ & $n$  & $k \equiv 1 \pmod{n}$ & $k \equiv 0\pmod{n}$ & Yes \\
$D_n$   & $2n$ & $n$  & $k \equiv 1 \pmod{n}$ & none                 & No  \\
$D_{nh}$& $4n$ & $n$  & $k \equiv 1 \pmod{n}$ & none                 & No  \\
$T$     & 12   & 12   & $k \equiv 1 \pmod{12}$& $k \equiv 0\pmod{12}$& Yes \\
$T_d$   & 24   & 12   & (reduced)             & (case-dep.)          & TBD \\
$O$     & 24   & 24   & $k \equiv 1 \pmod{24}$& $k \equiv 0\pmod{24}$& Yes \\
$O_h$   & 48   & 24   & (reduced)             & (case-dep.)          & TBD \\
$I$     & 60   & 60   & $k \equiv 1 \pmod{60}$& $k \equiv 0\pmod{60}$& Yes \\
$I_h$   & 120  & 60   & (reduced)             & (case-dep.)          & TBD \\
\bottomrule
\end{tabular}
\end{center}

% ============================================================
\section{Existence Theory}
% ============================================================

\subsection{Minimisation and the coercivity argument}

The action $\Action$ in \eqref{eq:single} is bounded below (the
potential is positive) but not coercive on all of $\mathcal{C}_N$
due to the possibility of collisions making $\Action \to +\infty$.
The key technique, introduced by Chenciner--Montgomery \cite{CM2000},
is to restrict minimisation to a homotopy class $[\gamma]_0 \in
\pi_0(\Lambda_N^G)$ that \emph{excludes collisions by topology}.

\begin{theorem}[Chenciner--Montgomery, 2000]\label{thm:cm}
If the homotopy class $[\gamma_0] \in \pi_1(\mathcal{C}_N / G)$
has the property that every path in this class must pass through the
collision set $\Delta$ to degenerate, then $\Action$ achieves its
infimum on $[\gamma_0] \cap \Fix(G)$, and the minimiser is a smooth
choreographic orbit.
\end{theorem}

The original proof uses the Marchal lemma \cite{Marchal2002}: any
orbit with a collision can be locally deformed to reduce the action
while remaining in the same homotopy class, ruling out collisions as
minimisers.

\subsection{Symmetry-forced collision exclusion}

For point groups $G$ with sufficiently high symmetry, collision
exclusion follows directly from the symmetry constraints without
invoking Marchal's lemma.

\begin{proposition}\label{prop:collision-free}
If the symmetry group $G$ contains elements $(R_k, k, +1)$ for all
$k = 1, \ldots, N-1$ such that $R_k$ has no fixed point in $\R^3$,
then the fixed-point subspace $\Fix(G)$ is collision-free.
\end{proposition}

\begin{proof}
At a collision between bodies $i$ and $j$, we have
$\bfq_i(t_0) = \bfq_j(t_0)$ at some time $t_0$.  By the
choreographic constraint, $\bfq_j(t) = \bfq_1(t + (j-i)/N)$, so
the collision becomes $\gamma(t_0) = \gamma(t_0 + (j-i)/N)$.  Under
the symmetry $(R_{j-i}, j-i, +1)$, this point must satisfy
$\gamma(t_0) = R_{j-i}\gamma(t_0)$, i.e.\ $\gamma(t_0)$ is a
fixed point of $R_{j-i}$.  If $R_{j-i}$ has no fixed point (e.g.\
if it is a rotation by an angle $\neq 0, \pi$ in $\SO(3)$),
no collision is possible.
\end{proof}

This is directly applicable to the $C_n$ and $C_{nv}$ families for
odd $n \geq 3$: the rotation $R_k = R_n^k$ has no fixed point for
$k \not\equiv 0 \pmod{n}$, so the odd-$n$ choreographies in these
families are automatically collision-free.

\subsection{Open existence questions}

The following cases are open and constitute part of the research
programme:
\begin{enumerate}[leftmargin=1.8em]
  \item Does the $C_{3v}$ class with $N=3$ admit a genuinely
        non-planar minimiser? (Numerically expected: yes.)
  \item Do any tetrahedral-class ($N=12$) minimisers exist in the
        $3D$-enabled modes?
  \item Does the icosahedral case ($N=60$) admit a minimiser, or
        does the small search space force the minimiser to collapse
        to a lower-symmetry orbit?
\end{enumerate}

% ============================================================
\section{Knot Theory of the Spacetime Trace}
% ============================================================

\subsection{The configuration space and its topology}

The unordered configuration space of $N$ distinct points in $\R^3$ is
\[
  \mathrm{UConf}_N(\R^3)
  \;=\;
  \bigl((\R^3)^N \setminus \Delta\bigr) / S_N.
\]
The fundamental group $\pi_1(\mathrm{UConf}_N(\R^3))$ is the braid
group $B_N$.  A choreographic orbit traces a loop in
$\mathrm{UConf}_N(\R^3)$ (since the bodies return to their starting
configuration after one period), and its homotopy class is a conjugacy
class in $B_N$ — the topological type of the orbit.

\subsection{From braid to knot: the spacetime trace}

A choreographic orbit defines a \emph{pure braid} in $B_N$: body 1
traces a closed curve and returns to its starting position (not
permuted to another position).  The spacetime trace of body 1 is the
closed curve
\begin{equation}\label{eq:strace}
  \Gamma : S^1 \to \R^3,
  \qquad
  \Gamma(\phi) = \gamma\!\left(\tfrac{\phi}{2\pi}\right),
\end{equation}
where we close the time-axis $[0,1]$ into $S^1$.  Topologically,
$\Gamma$ is a knot (or link) in $S^3 = \R^3 \cup \{\infty\}$.

For planar choreographies, $\Gamma$ lies in a solid torus
(thickening of the orbital plane times $S^1$) and is classified by
a braid word in $B_N$.  For three-dimensional choreographies,
$\Gamma$ can leave the solid torus and realise genuinely knotted
isotopy types not accessible from planar orbits.

\begin{remark}
The symmetry group $G$ constrains the knot type of $\Gamma$.  If
$(R, 0, +1) \in G$ is a spatial symmetry (no body permutation),
then $R$ acts on $\Gamma$ as a knot symmetry: $R \cdot \Gamma$ is
isotopic to $\Gamma$.  This places $\Gamma$ in the sub-collection of
knots admitting the symmetry $R$ — a significant restriction on
the knot table.
\end{remark}

\subsection{Polynomial invariants}

The following invariants are computed for each discovered orbit.

\textbf{Alexander polynomial.}  The Alexander polynomial
$\Delta_\Gamma(t) \in \Z[t^{\pm 1}]$ is computed from the Seifert
matrix $V$ of $\Gamma$ (a Seifert surface is triangulated from a
regular projection):
\[
  \Delta_\Gamma(t) = \det\!\left(t^{1/2}V - t^{-1/2}V^\top\right).
\]
The Alexander polynomial detects the unknot (it equals 1) and
distinguishes many knot types, but is not a complete invariant.

\textbf{HOMFLY-PT polynomial.}  The HOMFLY-PT polynomial
$P_\Gamma(v, z) \in \Z[v^{\pm 1}, z^{\pm 1}]$ satisfies the skein
relation
\[
  v^{-1}P_{L_+} - vP_{L_-} = z P_{L_0},
\]
where $L_+, L_-, L_0$ are diagrams differing at a single crossing.
It strictly dominates the Alexander polynomial (and the Jones
polynomial as a special case) and provides the finest easily-computable
invariant available.

\textbf{Hyperbolic volume.}  If the knot complement
$S^3 \setminus \Gamma$ admits a hyperbolic metric (the generic case by
Thurston's geometrisation theorem), its volume $\mathrm{vol}(\Gamma)$
is a complete knot invariant in the hyperbolic class.  Computed via SnapPy.

\subsection{The chemistry correspondence}

The class of \emph{torus knots} $T(p,q)$ (curves winding $p$ times
around the longitudinal direction and $q$ times around the meridional
direction of a torus) is of particular interest.  Torus knots are
the topological type of the most commonly synthesised molecular knots:
the trefoil $T(2,3)$ and the cinquefoil $T(2,5)$ have been realised
as single-molecule structures \cite{Stoddart2020}.

We conjecture that $C_n$-symmetric choreographies with $n = 2p+1$
preferentially realise torus knots of type $T(2, n)$: the
$z$-axis modes generate the ``tube'' direction of the torus while the
$xy$-plane modes generate the winding.  If confirmed, this establishes
a direct map from symmetry class to knot type and makes
the chemistry connection precise.

% ============================================================
\section{Floquet--Krein Stability Theory}
% ============================================================

\subsection{Linearisation}

Let $\mathbf{x}^*(t) = (\bfq^*(t), \dot\bfq^*(t)) \in T^*(\R^3)^N$
be a choreographic periodic orbit of period 1.  The linearised
equations of motion for a displacement $\delta\mathbf{x}(t)$ are
\begin{equation}\label{eq:linear}
  \frac{d}{dt}\delta\mathbf{x} \;=\; J(t)\,\delta\mathbf{x},
  \qquad
  J(t) = \begin{pmatrix} 0 & I \\ D^2U(\bfq^*(t)) & 0 \end{pmatrix},
\end{equation}
where $U = -\sum_{i<j}|\bfq_i - \bfq_j|^{-1}$ is the potential and
$D^2U$ is its Hessian matrix.  Equation \eqref{eq:linear} is a
linear Hamiltonian system with periodic coefficients.

\subsection{The monodromy matrix}

The fundamental matrix $\Phi(t)$ of \eqref{eq:linear} satisfies
$\dot\Phi = J(t)\Phi$, $\Phi(0) = I$.  The \emph{monodromy matrix}
$M = \Phi(1)$ encodes the linearised dynamics over one period.  Since
\eqref{eq:linear} is Hamiltonian, $M \in \mathrm{Sp}(6N, \R)$:
it is symplectic, $M^\top \Omega M = \Omega$, where $\Omega$ is the
standard symplectic form.

The eigenvalues of $M$ occur in groups consistent with the symplectic
structure:
\begin{itemize}[leftmargin=1.8em]
  \item Eigenvalues on the unit circle come in conjugate pairs
        $\{e^{i\theta}, e^{-i\theta}\}$ (elliptic modes).
  \item Real eigenvalues come in reciprocal pairs $\{\lambda,
        \lambda^{-1}\}$ (hyperbolic modes).
  \item Off-axis eigenvalues come in quadruples $\{\lambda,
        \bar\lambda, \lambda^{-1}, \bar\lambda^{-1}\}$
        (complex-hyperbolic modes).
\end{itemize}
By symmetry, $\lambda = 1$ always has multiplicity at least 4
(corresponding to time-translation and the three spatial
translations).  These are the \emph{trivial} eigenvalues.

\subsection{Krein signatures}

For elliptic eigenvalues $e^{\pm i\theta}$, linear stability is
not sufficient: one must also classify the Krein signature to
determine whether the orbit can be destabilised by arbitrarily small
perturbations.  The Krein signature of an elliptic eigenvalue
$\lambda = e^{i\theta}$ is the sign of the quadratic form
$\Omega(\cdot, M\cdot)$ restricted to the $\lambda$-eigenspace.

\begin{theorem}[Krein--Moser]\label{thm:km}
An elliptic eigenvalue $e^{i\theta}$ of $M$ is \emph{Krein-stable}
(i.e.\ cannot leave the unit circle under small perturbations of $M$
within $\mathrm{Sp}(6N,\R)$) if and only if its Krein signature is
definite (all $+1$ or all $-1$).  Instability requires the collision
of two eigenvalues with opposite Krein signatures.
\end{theorem}

\begin{corollary}
An orbit is \emph{Krein-stable} if all its non-trivial elliptic
eigenvalues have definite Krein signature.  Krein stability implies
linear stability and is a necessary condition for KAM non-linear
stability.
\end{corollary}

\subsection{KAM conditions}

By the KAM theorem applied to Hamiltonian systems near periodic
orbits \cite{Moser1962}, an orbit is non-linearly stable if:
\begin{enumerate}[leftmargin=1.8em]
  \item All non-trivial Floquet multipliers lie on the unit circle
        (linear stability).
  \item All Floquet angles $\theta_j$ satisfy non-resonance conditions
        $\sum_j n_j \theta_j \neq 2\pi k$ for all integer vectors
        $(n_j)$ with $\sum_j |n_j| \leq 4$ and integers $k$.
  \item A non-degeneracy (twist) condition holds for the
        Birkhoff normal form to fourth order.
\end{enumerate}
Conditions (1) and (2) are checkable from the Floquet spectrum;
condition (3) requires computing the fourth-order Birkhoff normal form,
which is numerically intensive but tractable for small $N$.

% ============================================================
\section{Implications for the Research Programme}
% ============================================================

The mathematical analysis above has several direct consequences for
the classification programme:

\begin{enumerate}[leftmargin=1.8em]
  \item \textbf{Non-planarity is a symmetry statement.}  Whether a
        given point group yields 3D orbits is determined exactly by
        Propositions~\ref{prop:cn}--\ref{prop:cnv}: the $z$-modes
        survive iff the group does not contain an element that
        eliminates them.  The summary table in Section~4 directly
        guides which groups to search.
  \item \textbf{Collision-free search is guaranteed for odd $C_n$
        and $C_{nv}$.}  Proposition~\ref{prop:collision-free} means
        the softening regularisation used in computation can be
        removed for these families without risk of spurious
        near-collision minimisers.
  \item \textbf{The knot type is constrained by the symmetry group.}
        The conjectured $T(2,n)$ torus knot structure for $C_n$
        choreographies, if confirmed numerically, would constitute
        a new theorem.  Its proof would use the symmetry-constrained
        form of $\Gamma$ and the classification of periodic braids
        under cyclic group actions.
  \item \textbf{The CompSpec is revised.}  The Reynolds projector
        construction in the CompSpec should now use the explicit
        mode selection rules derived here (Section~4) rather than the
        generic Reynolds averaging, which is equivalent but far more
        expensive to compute.  The mode selection table allows
        $P_G$ to be constructed in $O(K)$ rather than $O(|G| \cdot K)$
        time.
  \item \textbf{The RDD claim on icosahedral tractability is justified.}
        The ``dimensional miracle'' noted in Section~4.6 (the search
        space dimension at fixed $K/N$ is independent of $|G|$)
        formally justifies the optimistic treatment of icosahedral
        targets in the RDD.
\end{enumerate}

% ---------- REFERENCES ----------
\begin{thebibliography}{9}

\bibitem{CM2000}
A.~Chenciner and R.~Montgomery,
\textit{A remarkable periodic solution of the three-body problem in
the case of equal masses},
Ann.\ Math.\ \textbf{152} (2000), 881--901.

\bibitem{Marchal2002}
C.~Marchal,
\textit{How the method of minimization of action avoids singularities},
Cel.\ Mech.\ Dyn.\ Astr.\ \textbf{83} (2002), 325--353.

\bibitem{Palais1979}
R.~S.~Palais,
\textit{The principle of symmetric criticality},
Commun.\ Math.\ Phys.\ \textbf{69} (1979), 19--30.

\bibitem{Simo2002}
C.~Sim\'o,
\textit{New families of solutions in $N$-body problems},
in \textit{European Congress of Mathematics},
Birkh\"auser, Basel, 2001, pp.\ 101--115.

\bibitem{Moser1962}
J.~Moser,
\textit{On invariant curves of area-preserving mappings of an annulus},
Nachr.\ Akad.\ Wiss.\ G\"ottingen Math.-Phys.\ Kl.\ (1962), 1--20.

\bibitem{Stoddart2020}
F.~Schaufelberger et al.,
\textit{Synthesis of molecular knots},
Chem.\ Soc.\ Rev.\ \textbf{50} (2021), 3473--3500.

\end{thebibliography}

\end{document}
